\SourceLeaf{416}%
ergeben, jeweils dieselben Vorzeichen haben wie die Sinuswerte von \(\vartheta,\vartheta',\vartheta''\).

Ich darf auch nicht unerwähnt lassen, dass die angegebene allgemeine Lösung der Aufgabe in singulären Fällen entweder ihre Anwendbarkeit ganz verliert oder einer gewissen Anpassung bedarf. Ich beschränke mich hier jedoch darauf, die wichtigsten Punkte anzudeuten.

I. Ist einer der beobachteten Winkel \(=0\) oder \(=180^0\), so ist es vorteilhaft, den betreffenden Beobachtungsort als \((1)\) oder \((3)\) anzusetzen, auch wenn er zwischen den beiden anderen liegt. Wählt man Letzteres, bleiben alle Teile der allgemeinen Lösung gültig. Man muss lediglich \(b\) als unendlich groß betrachten und als Symbol in den Rechnungen beibehalten. In allen Ergebnissen fällt es dann von selbst weg. An mehr als einem Punkt darf der Winkel aber offenbar nicht \(0\) oder \(180^0\) sein, denn das kommt nur vor, wenn alle drei auf der Geraden \((4),(5)\) liegen. Dann wäre die Aufgabe unbestimmt.

II. Sind zwei der beobachteten Winkel gleich, so fallen zwei der Punkte \((1^*),(2^*),(3^*)\) zusammen, oder eine der Größen \(a,b,a+b\) wird \(=0\), also auch \(fF=0\). In diesem Fall geht die kubische Gleichung in eine quadratische über. Man sieht außerdem leicht, dass der erste Koeffizient der kubischen Gleichung \textit{nur} dann verschwindet, wenn zwei Winkel gleich sind.

III. Die angegebene Lösung ist nicht anwendbar, wenn die Größen \(A,B\) zu \(a,b\) proportional sind, also \(\lambda=0\) gilt. In diesem Fall wurde die kubische Gleichung zu Unrecht herangezogen. Sie enthält eine sachfremde Wurzel und die richtige Wurzel zweimal. Dann sind nämlich die beiden Kombinationen, durch die aus \([2]\) und \([3]\) die Gleichungen \([4]\) und \([5]\) abgeleitet wurden, offenkundig \textit{nicht verschieden}. Diese Gleichungen werden daher identisch, und jede für sich ergibt \(x=\frac A{2a}=\frac B{2b}\). Offenbar muss man dann aber noch eine der beiden Gleichungen \([2],[3]\) heranziehen. Ihre Kombination mit \([1]\) ergibt leicht \(cx=(t-n)\BareRoot{(1-xx)}\). Daraus folgt, dass die Gerade \(\mathrm{(I)}\) entweder durch den Punkt \((4)\) oder durch den Punkt \((5)\) geht, und
